Real-time machine learning (RTML) remains underutilised in neuroscience. This represents a missed opportunity—especially in the context of NaLoDuCo experimentation—where adaptive, low-latency computation could greatly enhance experimental control and scientific discovery.

We will attempt a real time implementation of the pipeline in Table 1, adapted for robustness in non-stationary evironments.

\subsubsubsection{Bonsai and Bonsai.ML}

Bonsai is a widely adopted open-source software ecosystem for experimental
control in neuroscience~\cite{lopesEtAl15}. With support from the
\href{https://gow.bbsrc.ukri.org/grants/AwardDetails.aspx?FundingReference=BB\%2FW019132\%2F1}{BBSRC},
we are developing infrastructure for intelligent, closed-loop experimentation
through the \href{https://bonsai-rx.org/machinelearning/}{Bonsai.ML} package.

We have already integrated several real-time ML models into Bonsai.ML,
including linear regression, linear dynamical systems, hidden Markov models,
and Bayesian point-process decoders. In collaboration with researchers at SWC
and UCL, we are applying these models to real-time inference of visual
receptive fields, foraging kinematics, behavioural state classification, and
spatial decoding of hippocampal spiking activity (see the
\href{https://bonsai-rx.org/machinelearning/examples/README.html}{Bonsai.ML examples}).

However, current Bonsai.ML models assume stationarity—a strong and often
invalid assumption in NaLoDuCo contexts. We will adapt these methods to operate
under non-stationary conditions using the previously described techniques.

Linear models from adaptive signal processing, such as least mean squares and
recursive least squares, are already suited for real-time use and require no
major modification. Nonlinear regression, particularly using artificial neural
networks (ANNs), can support real-time inference (especially on GPUs), but
real-time learning is more challenging. Continual learning methods—such as
experience replay and online regularization—introduce additional computational
overhead (e.g., data buffering, importance weight tracking) that may exceed the
constraints of real-time systems. To address these constraints, we will explore
adaptation strategies:

\begin{enumerate}

	\item\textbf{Real-time minibatch updates:} Training is performed on small
	minibatches as data arrive, balancing learning speed with latency.

	\item\textbf{Bounded-time updates:} Only a subset of model parameters
	(e.g., output layer weights) is updated to reduce computation time.

	\item\textbf{Hybrid strategy:} A dual-pipeline system in which a
	lightweight process performs fast inference, while a slower process
	(running on a separate processor or node) accumulates recent data in a
	rolling buffer and periodically retrains or updates the model. The updated
	model is deployed into the fast path only after validation, ensuring
	continuity in inference.

\end{enumerate}

All new RTML tools will be released as open-source extensions to the
\href{https://bonsai-rx.org/machinelearning/}{Bonsai.ML} package.

\subsubsubsection{Experimental Replays in AEON}

We will extend the AEON platform with functionality to replay previously
recorded experimental data. Real-time machine learning (RTML) objects in AEON
typically receive input from live data streams—such as video frames from
cameras or voltage measurements from Neuropixels probes—each precisely
timestamped using the HARP protocol. This protocol ensures that every item in
the stream is tagged with the exact time it occurred during the experiment.

The new replay feature will regenerate all data streams from a previous
experiment and deliver each data item to its observers at exactly the same
relative time as in the original experiment. This enables RTML components to be
tested under identical conditions to those in live experiments, while
eliminating the need to rerun time-consuming and resource-intensive
experiments. This supports the 3Rs by reducing animal use (Reduction) and
enabling partial replacement of live testing (Replacement).

Replay functionality will enable comprehensive testing of RTML methods.
Researchers will be able to debug their pipelines, evaluate different parameter
settings, and verify that real-time constraints are met—all without placing
additional demands on experimental animals or facilities. This is particularly
critical in NaLoDuCo settings, where experiments run continuously for weeks or
months and are costly, laborious, and ethically sensitive.

\subsubsubsection{Deliverables}

\begin{enumerate}

	\item New real-time machine learning methods for non-stationary
	experimental control, integrated into the open-source
	\href{https://bonsai-rx.org/machinelearning/}{Bonsai.ML} package.

	\item AEON experimental replay functionality enabling real-time
	reproductions of previously recorded AEON experiments, saving time,
	resources and animal experiments.

	\item Peer-reviewed publications co-authored with SWC and AIND researchers,
	showcasing scientific advances made possible by the new RTML capabilities.

\end{enumerate}
